\subsection{模的直和}

设 \((\mathbf{E}_t)_{t \in I}\) 是一族 A-模， \(\mathbf{F} = \prod_{i \in I} \mathbf{E}_i\) 是它们的积。集合 \(\mathbf{E}\) 由 \(x \in F\) 中使 \(\operatorname{pr}_t x = 0\) 除有限个指标外均成立的元素组成，它显然是 \(\mathbf{F}\) 的一个子模，称为模族 \((\mathbf{E}_t)\) 的外直和（或简称直和），记作 \(\bigoplus_{i \in I} \mathbf{E}_i\) （I，§4，第 9 号）。当 I 有限时，有 \(\bigoplus_{i \in I} \mathbf{E}_i = \prod_{i \in I} \mathbf{E}_t\) ；若 \(\mathbf{I} = [p, q]\) （ \(\mathbf{Z}\) 的区间），我们也记

\[
\bigoplus_ {i \in I} \mathbf {E} _ {i} = \mathbf {E} _ {p} \oplus \mathbf {E} _ {p + 1} \oplus \dots \oplus \mathbf {E} _ {q}.
\]

对所有 \(\kappa \in I\) ，设 \(j_{\kappa}\) 是映射 \(\mathbf{E}_{\kappa} \to F\) ，它把每个 \(x_{\kappa} \in \mathbf{E}_{\kappa}\) 对应到 \(F\) 中满足 \(\operatorname{pr}_{\iota}(j_{\kappa}(x_{\kappa})) = 0\) （对 \(\iota \neq \kappa\) ）以及 \(\operatorname{pr}_{\kappa}(j_{\kappa}(x_{\kappa})) = x_{\kappa}\) 的元素；立即可以看出， \(j_{\kappa}\) 是 \(\mathbf{E}_{\kappa}\) 到直和 \(E\) （诸 \(\mathbf{E}_{\iota}\) 的直和）中的单射线性映射，我们称之为典范单射；子模 \(j_{\kappa}(\mathbf{E}_{\kappa})\) （属于 \(E\) ，且同构于 \(\mathbf{E}_{\kappa}\) ）称为 \(E\) 的指标为 \(\kappa\) 的分量子模。通常将它与 \(\mathbf{E}_{\kappa}\) 等同（借助 \(j_{\kappa}\) ）。

因此，对所有 \(x \in \mathbf{E} = \bigoplus_{i \in I} \mathbf{E}_i\) ，我们有

\[
x = \sum_ {i \in I} j _ {i} (\mathrm{pr} _ {i} x).\tag{20}
\]

命题 6. 设 \((\mathbf{E}_t)_{t \in I}\) 是一族 A-模，M 是一个 A-模，并且对所有 \(t \in I\) ，设 \(f_t: \mathbf{E}_t \to \mathbf{M}\) 是线性映射。则存在且仅存在一个线性映射 \(g: \bigoplus_{t \in I} \mathbf{E}_t \to \mathbf{M}\) ，使得对所有 \(t \in I\) ：

\[
g \circ j _ {i} = f _ {i}.\tag{21}
\]

根据 (20)，若 \(g\) 存在，则必然对所有 \(x \in \bigoplus_{\iota \in I} \mathbf{E}_{\iota}\) 有

\[
g (x) = \sum_ {i} g (j _ {i} (\mathrm{pr} _ {i} (x))) = \sum_ {i} f _ {i} (\mathrm{pr} _ {i} (x)),
\]

由此得到 \(g\) 的唯一性。反之，令 \(g(x) = \sum_{i} f_i(\mathrm{pr}_i(x))\) （对所有 \(x \in \bigoplus_{i \in I} E_i\) ），可以立即验证这样就定义了一个满足命题条件的线性映射。

当不致引起混淆时，我们记 \(g = \sum_{i \in I} f_i\) （当族 \((f_i)\) 不是有限支集时，这与 I，§2，第 1 号的约定相反）。

特别地，若 J 是 I 的任意子集，则典范单射 \(j_{i}\) （ \(i \in J\) ）定义了一个典范线性映射 \(j_{J}: \bigoplus_{i \in J} E_{i} \to \bigoplus_{i \in I} E_{i}\) ，它把每个 \((x_{i})_{i \in I}\) 对应到满足下列条件的元素 \((x_{i}')_{i \in I}\) ： \(x_{i}' = x_{i}\) （对 \(i \in J\) ）， \(x_{i}' = 0\) （对 \(i \notin J\) ）；这个映射显然是单射。此外，若 \((J_{\lambda})_{\lambda \in L}\) 是 I 的一个划分，则映射 \(i: \bigoplus_{\lambda \in L} \left( \bigoplus_{i \in J_{\lambda}} E_{i} \right) \to \bigoplus_{i \in I} E_{i}\) （由命题 6，它对应于族 \((j_{J_{\lambda}})\) ）是一个同构，称为典范同构（直和的“结合性”）。

推论 1. 设 \((\mathbf{E}_{\iota})_{\iota \in I}, (\mathbf{F}_{\lambda})_{\lambda \in L}\) 是两个 A-模族。映射

\[
\operatorname{Hom} _ {\mathbf {A}} \left(\bigoplus_ {t \in I} \mathrm{E} _ {t}, \prod_ {\lambda \in L} \mathrm{F} _ {\lambda}\right)\rightarrow \prod_ {(t, \lambda) \in I \times L} \operatorname{Hom} _ {\mathbf {A}} \left(\mathrm{E} _ {t}, \mathrm{F} _ {\lambda}\right)\tag{22}
\]

它把每个 \(g \in \mathrm{Hom}_{\mathbf{A}} \left( \bigoplus_{t \in I} \mathbf{E}_t, \prod_{\lambda \in L} \mathbf{F}_\lambda \right)\) 对应到族 \((\mathbf{pr}_\lambda \circ g \circ j_t)\) ，这是一个 Z-模同构（称为典范同构）。

这由命题 6 和第 5 号的命题 4 得出。

推论 2. 设 \((\mathbf{E}_t)_{t \in I}\) 是一族 A-模， \(F\) 是一个 A-模，并且对每个 \(t \in I\) ，设 \(f_t: \mathbf{E}_t \to F\) 是线性映射。要使 \(f = \sum_{t \in I} f_t\) 是 \(\mathbf{E} = \bigoplus_{t \in I} \mathbf{E}_t\) 到 \(F\) 上的同构，必须且只需对每个 \(t \in I\) 存在具有下列性质的线性映射 \(g_t: \mathbf{F} \to \mathbf{E}_t\) ：

\begin{enumerate}
\def\labelenumi{(\arabic{enumi})}
\item
  \(g_{i} \circ f_{i} = 1_{E_{i}}\) 对所有 \(i \in I\) 成立。
\item
  \(g_{\iota} \circ f_{\kappa} = 0\) 对 \(\iota \neq \kappa\) 。
\item
  对所有 \(y \in \mathbf{F}\) ，族 \((g_{\iota}(y))\) 具有有限支集，且
\end{enumerate}

\[
y = \sum_ {i \in I} f _ {i} (g _ {i} (y)).\tag{23}
\]

注意，若 I 有限，则最后一个条件也可写为

\[
\sum_ {i \in I} f _ {i} \circ g _ {i} = 1 _ {\mathbf {F}}.\tag{24}
\]

显然这些条件是必要的，因为它们被 \(g_{\iota} = \mathrm{pr}_{\iota} \circ f\) 满足。反之，若它们成立，则对所有 \(y \in \mathbf{F}\) ， \(g(y) = \sum_{\iota} j_{\iota}(g_{\iota}(y))\) 有定义，并且立即可以看出 \(g\) 是 \(\mathbf{F}\) 到 \(\mathbf{E}\) 中的线性映射。对所有 \(y \in \mathbf{F}\) ，根据假设有 \(f(g(y)) = \sum_{\iota \in \mathbf{I}} f_{\iota}(g_{\iota}(y)) = y\) 。另一方面，对所有 \(x \in \mathbf{E}\) ，根据假设有 \(g_{\kappa}(f(x)) = g_{\kappa}\left(\sum_{\iota} f_{\iota}(\mathrm{pr}_{\iota}(x))\right) = g_{\kappa}(f_{\kappa}(\mathrm{pr}_{\kappa}(x))) = \mathrm{pr}_{\kappa}(x)\) ；

因此 \(g(f(x)) = \sum_{i} j_i(g_i(f(x))) = \sum_{i} j_i(\mathrm{pr}_i(x)) = x\) ，这就证明了该推论。

命题 7. (i) 设 \((\mathbf{E}_t)_{t \in I}\) , \((\mathbf{F}_t)_{t \in I}\) 是具有同一指标集 I 的两个 A-模族；对每族线性映射 \(f_t: \mathbf{E}_t \to \mathbf{F}_t\) ( \(t \in I\) )，在 \(\bigoplus_{t \in I} \mathbf{E}_t\) 上，线性映射 \((x_t) \mapsto (f_t(x_t))\) 的限制是一个线性映射 \(f: \bigoplus_{t \in I} \mathbf{E}_t \to \bigoplus_{t \in I} \mathbf{F}_t\) ，记作 \(\bigoplus_{t \in I} f_t\) 或 \(\bigoplus_{t} f_t\) （其中 \(f = f_p \oplus f_{p+1} \oplus \cdots \oplus f_q\) ，若 \(\mathbf{I} = [p, q]\) 是 \(\mathbf{Z}\) 的区间）.

\begin{enumerate}
\def\labelenumi{(\roman{enumi})}
\setcounter{enumi}{1}
\tightlist
\item
  设 \((\mathbf{G}_{\iota})_{\iota \in \mathbf{I}}\) 是以 \(\mathbf{I}\) 为指标集的第三族 A-模，并且对所有 \(\iota \in \mathbf{I}\) ，设 \(g_{\iota}: F_{\iota} \to G_{\iota}\) 是线性映射；我们记 \(g = \bigoplus_{\iota} g_{\iota}\) 。要使每个序列 \(\mathbf{E}_{\iota} \xrightarrow{f_{\iota}} F_{\iota} \xrightarrow{g_{\iota}} G_{\iota}\) 都正合，必须且只需序列
\end{enumerate}

\[
\bigoplus_ {i} \mathrm{E} _ {i} \xrightarrow {f} \bigoplus_ {i} \mathrm{F} _ {i} \xrightarrow {g} \bigoplus_ {i} \mathrm{G} _ {i}
\]

是正合的。

显然，对所有 \((x_{i})\in \bigoplus_{i}\mathbf{E}_{i}\) ，族 \((f_{i}(x_{i}))\) 具有有限支集，由此得 (i)。另一方面，说元素 \(y = (y_{i})\) （ \(\bigoplus_{i}\mathbf{F}_{i}\) 中的元素）属于 \(\operatorname {Ker}(g)\) ，就是说 \(y_{i}\in \operatorname {Ker}(g_{i})\) 对所有 \(i \in I\) 成立（第 5 号，命题 5）；类似地，若 \(y_{i}\in \operatorname {Im}(f_{i})\) 对所有 \(\iota \in I\) 成立，则对每个 \(\iota \in I\) 存在 \(x_{i}\in \mathbf{E}_{i}\) 使得 \(y_{i} = f_{i}(x_{i})\) ，并且当 \(y_{i} = 0\) 时可以假设 \(x_{i} = 0\) ；因此 \(y\in \operatorname {Im}(f)\) ，而逆命题是显然的。

推论 1. 在命题 7 (i) 的条件下，

\[
\operatorname{Ker} (f) = \bigoplus_ {i \in I} \operatorname{Ker} (f _ {i}), \quad \operatorname{Im} (f) = \bigoplus_ {i \in I} \operatorname{Im} (f _ {i})\tag{25}
\]

并且存在典范同构

\[
\operatorname{Coim} (f) \cong \bigoplus_ {i \in I} \operatorname{Coim} (f _ {i}), \quad \operatorname{Coker} (f) \cong \bigoplus_ {i \in I} \operatorname{Coker} (f _ {i})
\]

其定义方式与第 5 号命题 5 的推论中相同。特别地，要使 \(f\) 是单射（相应地，满射、双射、零映射），必须且只需每个 \(f_i\) 都是单射（相应地，满射、双射、零映射）。

若对所有 \(t \in I\) ，我们考虑一个子模 \(F_t\) （ \(E_t\) 的子模），则模 \(\bigoplus_{t \in I} F_t\) 是 \(\bigoplus_{t \in I} E_t\) 的一个子模，并且根据命题 7 的推论 1，存在典范同构

\[
\bigoplus_ {i \in I} \left(\mathrm{E} _ {i} / \mathrm{F} _ {i}\right)\cong \left(\bigoplus_ {i \in I} \mathrm{E} _ {i}\right) / \left(\bigoplus_ {i \in I} \mathrm{F} _ {i}\right).\tag{26}
\]

推论 2. 设 \((\mathbf{E}_t)_{i \in I}, (\mathbf{E}_t')_{i \in I}, (\mathbf{F}_\lambda)_{\lambda \in L}, (\mathbf{F}_\lambda')_{\lambda \in L}\) 是四个 A-模族，并且对每个 \(i \in I\) （相应地，每个 \(\lambda \in L\) ）， \(u_i: \mathbf{E}_i' \to \mathbf{E}_i\) （相应地 \(v_\lambda: \mathbf{F}_\lambda \to \mathbf{F}_\lambda'\) ）是线性映射。则图表

\[
\begin{array}{c} \operatorname{Hom} \left(\bigoplus_ {t \in I} E _ {t} ^ {\prime}, \prod_ {\lambda \in L} F _ {\lambda} ^ {\prime}\right) \xrightarrow {\phi^ {\prime}} \prod_ {(t, \lambda) \in I \times L} \operatorname{Hom} (E _ {t} ^ {\prime}, F _ {\lambda} ^ {\prime}) \\ \operatorname{Hom} \left(\bigoplus_ {t} u _ {t}, \prod_ {\lambda} v _ {\lambda}\right) \Bigg | ^ {\uparrow} \\ \operatorname{Hom} \left(\bigoplus_ {t \in I} E, \prod_ {\lambda \in L} F _ {\lambda}\right) \xrightarrow [ \phi ]{\quad} \prod_ {(t, \lambda) \in I \times L} \operatorname{Hom} (E _ {t}, F _ {\lambda}) \end{array}
\]

（其中 \(\phi\) 和 \(\phi'\) 是命题 6 的推论 1 中定义的典范同构）是交换的。

验证由定义立即得出。

当所有 \(E_{t}\) 都等于同一个 A-模 E 时，直和 \(\bigoplus_{i\in I}E_{i}\) 也记作 \(\mathrm{E}^{(I)}\) ：其元素是 I 到 E 中的具有有限支集的映射。若对所有 \(t,f_{t}\) 都取为恒等映射 \(E\to E\) ，则由命题 6 得到一个线性映射 \(\mathrm{E}^{(I)}\to\mathrm{E}\) ，称为余对角映射，它把 E 中元素的每个有限支集族 \((x_{t})_{t\in I}\) 对应到其和 \(\sum_{i\in I}x_{i}\) 。

注记。回顾，直和的定义可立即推广到一族未必交换的群 \((\mathrm{E}_{\iota})_{\iota\in\mathbb{I}}\) ，当然以乘法记号代替加法记号；此时我们说“限制积”或“限制和”以代替“直和”（I，§4，第 9 号）。

注意，E 是乘积 \(F = \prod_{i \in I} E_{i}\) 的一个正规子群，并且每个 \(j_{\kappa}(E_{\kappa})\) 都是 F 的正规子群；此外，对两个不同的指标 \(\lambda, \mu\) ， \(j_{\lambda}(E_{\lambda})\) 的每个元素与 \(j_{\mu}(E_{\mu})\) 的每个元素可交换。命题 6 可推广到一般的情形，但需假设：对两个不同的指标 \(\lambda, \mu\) ， \(f_{\lambda}(E_{\lambda})\) 的每个元素在 M 中与 \(f_{\mu}(E_{\mu})\) 的每个元素可交换（I，§4，第 9 号，命题 12）。限制和的“结合性”性质由此立即得出。命题 7 及其推论 1 和 2 无需修改即成立。

